% SOURCE_AUTHORITY: sga5_fr_workpass.tex, lines 7033--7057 (LF-normalized unit).
% SOURCE_SHA256: 791F4EFFC5E02832D5D77ED03518C8156D6F07E4C8238B03545DB93D883FBB28
\begin{statement}{Proposição 2.4.4}
\begin{enumerate}[label=(\roman*)]
\item A categoria $P_{AR}$ é abeliana.
\item O funtor $p_{AR}:P\to P_{AR}$ é exato.
\item Para que um objeto $X$ de $P$ seja AR-nulo, é necessário e suficiente que $p_{AR}(X)=0$.
\item O par $(p_{AR},P_{AR})$ é um quociente da categoria abeliana $P$ pela subcategoria abeliana espessa dos objetos AR-nulos.
\end{enumerate}
\end{statement}

\begin{prooftext}[Demonstração]
Se for admitido (iii), as outras afirmações são consequências de (2.3.5). De (2.3.4) deduz-se que, para que um objeto $X$ de $P$ seja AR-nulo, é necessário e suficiente que $p_{AR}(\id_X)=0$, de onde segue-se (iii).
\end{prooftext}

\begin{statement}{Corolário 2.4.5}
Seja $u:X\to Y$ da categoria $P$. As seguintes afirmações são equivalentes:
\begin{enumerate}[label=(\roman*)]
\item $p_{AR}(u)$ é um isomorfismo.
\item O núcleo e o conúcleo de $u$ são AR-nulos.
\end{enumerate}
Um morfismo $u:X\to Y$ da categoria $P$ dos sistemas projetivos que verifique as condições equivalentes (i) e (ii) será chamado AR-isomorfismo.
\end{statement}

\begin{statement}{Definição 2.4.6}
Dois objetos $X$ e $Y$ da categoria $P$ são chamados de AR-isomorfos se $p_{AR}(X)$ e $p_{AR}(Y)$ são isomorfos na categoria $P_{AR}$. Se $X$ e $Y$ são dois objetos AR-isomorfos da categoria $P$, não existe geralmente um AR-isomorfismo $u$ de $X$ em $Y$.
\end{statement}
